%%═══════════════════════════════════════════════════════════════════
%% Lecture 05 — Gradient Descent, Newton's Optimization, and
%%              Introduction to ODEs with Euler's Method
%% 77315 — Computational Physics
%% Date: 30/04/2026
%%═══════════════════════════════════════════════════════════════════

\section{Lecture 05 — Gradient Descent, Newton's Optimization Method, and
  Introduction to ODEs}\label{sec:lec05}

\begin{center}
\begin{tabular}{ll}
  \textbf{Date:}\quad 30/04/2026 & \\
\end{tabular}
\end{center}
\hrule\vspace{1.5em}

%%──────────────────────────────────────────────────────────────────
\subsection*{Overview}
This lecture transitions from root-finding (Lecture~\ref{sec:lec04}) to
\emph{optimization} — finding a point where the gradient vanishes — and then
opens the topic of \emph{numerical integration of ordinary differential
equations}. The two topics are related: the gradient-descent step is itself
governed by a continuous-time ODE, and the Euler method we derive here is the
simplest possible discretisation of any such ODE.

%%──────────────────────────────────────────────────────────────────
\subsection{Gradient Descent}

\noindent\textbf{Physical intuition:} Imagine a blindfolded hiker on a
multidimensional landscape trying to reach the lowest valley. At each step
the hiker measures the slope in every direction (the gradient $\grad f$) and
steps in the \emph{opposite} direction — directly downhill. This is gradient
descent.

\subsubsection{Derivation via Taylor Expansion}
Let $f : \mathbb{R}^N \to \mathbb{R}$ be the objective function. Expand to
first order around the current point $\mathbf{x}_0$:
\begin{equation}
f(\mathbf{x}_0 + \delta) \approx f(\mathbf{x}_0) + (\grad f)\cdot\delta
\end{equation}
To decrease $f$ as rapidly as possible, choose $\delta$ antiparallel to the
gradient:
\begin{equation}
\boxed{\delta = -\alpha\,\grad f(\mathbf{x}_0)}
\end{equation}
where $\alpha > 0$ is the \textbf{step size} (learning rate). The update rule
is therefore:
\begin{equation}
\mathbf{x}_{n+1} = \mathbf{x}_n - \alpha\,\grad f(\mathbf{x}_n)
\end{equation}

\subsubsection{Step Size Selection via 1D Line Search}
Choosing $\alpha$ is non-trivial:
\begin{itemize}
    \item Too large: the update overshoots the minimum and $f$ may increase.
    \item Too small: convergence is extremely slow.
\end{itemize}
\noindent\textbf{Optimal strategy:} treat $\alpha$ as a 1D optimisation
problem. Define $g(\alpha) = f(\mathbf{x}_n - \alpha\,\grad f)$ and minimise
it using Golden Section Search (Lecture~\ref{sec:lec04},
\S\ref{subsec:linesearch}). This \emph{exact line search} is expensive but
guarantees the best step along the descent direction. In practice, the
\emph{backtracking} (Armijo) criterion from Lecture~\ref{sec:lec04} is used
instead — it is much cheaper and still ensures strict descent.

\subsubsection{Convergence Properties and Limitations}
\begin{itemize}
    \item Gradient descent converges to a \textbf{local} minimum, not
    necessarily the global one. The algorithm has no mechanism to escape a
    local basin.
    \item Near the minimum, convergence is \emph{linear} (error decreases by
    a constant factor per step). In contrast to Newton's method (below),
    which is quadratic, gradient descent can require many more steps in
    poorly scaled or ill-conditioned problems.
    \item The method is inefficient in narrow, elongated valleys (ravines):
    the gradient oscillates across the valley rather than proceeding along it.
\end{itemize}

%%──────────────────────────────────────────────────────────────────
\subsection{Newton's Method for Optimization (Second-Order Methods)}

\noindent\textbf{Physical intuition:} Gradient descent uses only the slope
(first derivative) of the landscape. Newton's method adds the \emph{curvature}
(second derivative) — it fits a local paraboloid to the landscape and jumps
directly to its bottom, rather than crawling downhill.

\subsubsection{Second-Order Taylor Expansion and the Hessian}
Expand $f$ to second order around $\mathbf{x}_0$:
\begin{equation}
f(\mathbf{x}_0 + \delta)
\approx f(\mathbf{x}_0) + (\grad f)\cdot\delta
+ \tfrac{1}{2}\,\delta^T \Hess\,\delta
\end{equation}
where the \textbf{Hessian matrix} $\Hess$ has entries
\begin{equation}
H_{ij} = \frac{\partial^2 f}{\partial x_i\,\partial x_j}
\end{equation}
The local paraboloid has a unique minimum where its gradient vanishes:
\begin{equation}
\grad_\delta\bigl[(\grad f)\cdot\delta + \tfrac{1}{2}\,\delta^T\Hess\,\delta\bigr]
= \grad f + \Hess\,\delta = 0
\end{equation}
This gives the \textbf{Newton step}:
\begin{equation}
\boxed{\Hess\,\delta = -\grad f}
\end{equation}

\subsubsection{Solving the Newton System}
The equation $\Hess\,\delta = -\grad f$ is a square linear system of size
$N\times N$, solved by Gaussian elimination in $\bigO{N^3}$ operations —
exactly the same infrastructure as multidimensional Newton-Raphson for roots
(Lecture~\ref{sec:lec04}).

The new iterate is $\mathbf{x}_{\text{new}} = \mathbf{x} + \delta$. Near
the minimum, convergence is \emph{quadratic}: the number of correct digits
roughly doubles each iteration.

\subsubsection{Positive Definiteness Condition}
For the paraboloid to have a minimum (not a maximum or saddle point), the
Hessian must be \textbf{Positive Definite} (PD): all eigenvalues
$\lambda_i > 0$, equivalently $\mathbf{v}^T\Hess\,\mathbf{v} > 0$ for all
non-zero $\mathbf{v}$.

If $\Hess$ is not PD at the current point, the Newton step points toward a
saddle or maximum — we are not near a minimum. In this case, one falls back
to gradient descent or adds a positive multiple of the identity matrix
(\emph{Levenberg-Marquardt regularisation}) to make $\Hess + \mu I$ PD.

%%──────────────────────────────────────────────────────────────────
\subsection{Checking Positive Definiteness via Gaussian Elimination}

\noindent\textbf{Physical intuition:} Computing all $N$ eigenvalues of $\Hess$
to verify positivity costs $\bigO{N^3}$ with a large prefactor. Gaussian
elimination on a symmetric matrix offers an equally conclusive check at lower
cost.

\subsubsection{The Criterion}
Apply Gaussian elimination (without row swaps, since $\Hess$ is symmetric)
to $\Hess$. At each step, divide the $k$-th row by the diagonal pivot
$B_{kk}$. The key theorem is:

\begin{tcolorbox}[colback=blue!5!white,colframe=blue!75!black,
  title=Positive Definiteness via Gaussian Elimination]
A real symmetric matrix $\Hess$ is Positive Definite \emph{if and only if}
all diagonal pivot elements $B_{kk}$ remain strictly positive during
Gaussian elimination without pivoting.

\textbf{Why:} The pivots $B_{11}, B_{22}, \ldots, B_{NN}$ are the ratios of
successive leading principal minors:
\[
B_{kk} = \frac{\det(\Hess_k)}{\det(\Hess_{k-1})}
\]
By Sylvester's criterion, $\Hess$ is PD $\iff$ all leading principal minors
are positive $\iff$ all pivots are positive. If any $B_{kk} \leq 0$, the
matrix is not PD and we are not at a (local) minimum.
\end{tcolorbox}

\subsubsection{Connection to the Minimization Check (Lecture~\ref{sec:lec04})}
In Lecture~\ref{sec:lec04} we noted that checking positive definiteness of
$\Hess$ via Gaussian elimination is cheaper than computing eigenvalues. The
current lecture provides the algorithmic detail: run standard Gaussian
elimination on $\Hess$ (which is symmetric, so only the upper triangle need
be stored) and inspect the diagonal pivots. Cost: $\bigO{N^3/6}$, a factor of
6 cheaper than full LU with the same asymptotic scaling.

%%──────────────────────────────────────────────────────────────────
\subsection{Ordinary Differential Equations: Problem Setup}

\noindent\textbf{Physical intuition:} Almost all of classical physics is
expressed as ODEs: Newton's second law, circuit equations, orbital mechanics.
The goal is to integrate these equations forward in time to find the
trajectory $y(x)$ given an initial condition.

\subsubsection{Standard Form: First-Order System}
The general first-order ODE is:
\begin{equation}
\frac{dy}{dx} = f\bigl(x,\, y(x)\bigr), \qquad y(x_0) = y_0
\end{equation}
Higher-order equations are reduced to this form by introducing auxiliary
variables. Consider Newton's second law:
\begin{equation}
m\,\ddot{r} = F(t, r, \dot{r})
\end{equation}
Define $v = \dot{r}$. This becomes a \emph{system} of two first-order ODEs:
\begin{align}
\dot{r} &= v \\
\dot{v} &= \frac{F(t, r, v)}{m}
\end{align}
In vector form $\mathbf{y} = (r, v)^T$:
\begin{equation}
\frac{d\mathbf{y}}{dt} = \mathbf{f}(t, \mathbf{y})
\end{equation}
All standard numerical ODE solvers work on this canonical first-order system
form, so the reduction is not merely a notational convenience — it is the
required input format.

%%──────────────────────────────────────────────────────────────────
\subsection{Euler's Method}

\noindent\textbf{Physical intuition:} Euler's method is the simplest possible
numerical ODE integrator. At each step, assume the derivative $f(x_n, y_n)$
is constant over the interval $[x_n, x_{n+1}]$ and use it to extrapolate
linearly.

\subsubsection{Derivation}
Discretise the domain into uniform steps of size $h$:
$x_n = x_0 + n\,h$. Replace the derivative by its forward-difference
approximation:
\begin{equation}
\frac{y_{n+1} - y_n}{h} \approx f(x_n, y_n)
\end{equation}
Solving for $y_{n+1}$:
\begin{equation}
\boxed{y_{n+1} = y_n + h\,f(x_n,\, y_n)}
\end{equation}

\subsubsection{Error Analysis}
\textbf{Local (per-step) error.} Expand the exact solution in Taylor series:
\begin{equation}
y(x_{n+1}) = y(x_n) + h\,y'(x_n) + \frac{h^2}{2}\,y''(x_n) + \bigO{h^3}
\end{equation}
Since $y'(x_n) = f(x_n, y_n)$, the Euler update matches the first two terms.
The local truncation error is therefore:
\begin{equation}
e_{\text{local}} = \frac{h^2}{2}\,y''(x_n) + \bigO{h^3} = \bigO{h^2}
\end{equation}

\textbf{Global error.} Over the full domain $[x_0, X]$ there are $n =
(X - x_0)/h$ steps, so errors accumulate:
\begin{equation}
e_{\text{global}} \approx n \cdot e_{\text{local}} = \frac{X - x_0}{h}
\cdot \bigO{h^2} = \bigO{h}
\end{equation}
Euler's method is therefore a \textbf{first-order} method: halving the step
size halves the global error.

\begin{tcolorbox}[colback=red!5!white,colframe=red!75!black,
  title=Why Euler's Method is Rarely Used in Practice]
Although simple and instructive, Euler's method suffers from two problems:
\begin{enumerate}
    \item \textbf{Accuracy}: Only first-order global accuracy. To achieve
    $10^{-6}$ absolute error over unit interval requires $h \sim 10^{-6}$,
    i.e., one million evaluations of $f$.
    \item \textbf{Stability}: For stiff equations (where $f$ has a large
    Lipschitz constant), the step size must satisfy $h < 2/|f_y|$ or the
    error grows exponentially. This can force $h$ to be astronomically small
    even when the solution itself is smooth.
\end{enumerate}
Higher-order methods (Runge-Kutta, Adams-Bashforth) achieve much better
accuracy per function evaluation. The Runge-Kutta 4 (RK4) method, covered in
the next lecture, achieves $\bigO{h^4}$ local error ($\bigO{h^4}$ global)
with only 4 evaluations of $f$ per step.
\end{tcolorbox}

%%──────────────────────────────────────────────────────────────────
\subsection*{Recommended Reading}
\begin{itemize}
    \item \textit{Numerical Recipes}, W.H. Press et al.\ --- Ch.\ 10.6
    (Gradient Methods; Conjugate Gradients), Ch.\ 16.1 (Runge-Kutta Method;
    Euler as special case).
    \item \textit{Computational Physics}, Mark Newman --- Ch.\ 8 (Ordinary
    Differential Equations).
    \item \textit{An Introduction to Computational Physics}, Tao Pang ---
    Ch.\ 3 (Ordinary Differential Equations).
\end{itemize}
